\section{Epistemic status of the construction}

To avoid attributing our synthesis to the source literature, Table~\ref{tab:status} separates
established results from the present construction.

\begin{table}[ht]
\centering
\small
\begin{tabular}{p{0.30\linewidth}p{0.20\linewidth}p{0.39\linewidth}}
\toprule
Statement & Status & Basis\\
\midrule
Costly knowledge acquisition and communication can induce knowledge hierarchies and
management by exception.
& Established
& Garicano (2000).\\[2mm]

Productive experience can generate learning.
& Established
& Learning-by-doing literature; Arrow (1962).\\[2mm]

Learning can be task-specific.
& Established
& Task-specific human-capital literature; Gibbons and Waldman (2004).\\[2mm]

Organizational knowledge can change as a function of task-performance experience; knowledge
is created, retained and transferred.
& Established framework
& Argote and Miron-Spektor (2011).\\[2mm]

Work allocation can explicitly affect future competence through learning and depreciation.
& Established formal model
& Gutjahr (2011), under its assumptions.\\[2mm]

Allocation can determine real work, which may generate task-dependent learning exposure
entering a competence transition.
& Candidate / hypothesis
& Requires explicit \(\mathcal W\), \(\mathcal E\), and \(\mathcal L\) mappings in the HLS ontology.\\[2mm]

\((F_t,S_t)\to a_t\to w_t\to e_t\to S_{t+1}\to a_{t+1}\).
& Candidate HLS scaffold
& Not stated as such by any single cited source.\\[2mm]

A particular machine-learning implementation of \(S_t,\mathcal R,\mathcal W,\mathcal E,\mathcal L\).
& Open
& Requires later modeling and empirical work.\\
\bottomrule
\end{tabular}
\caption{Epistemic status of the two-beam skeleton.}
\label{tab:status}
\end{table}

The individual ingredients are source-grounded. Their mapping and combination should be
described as a \emph{candidate modeling scaffold}, not as an established theorem or a
completed transfer inherited from the literature.
